\chapter{2015年湖北卷（理）}

\section{单选题}

\begin{problem}
\(\i\) 为虚数单位，\(\i^{607}\) 的共轭复数为（\quad）

\choices
    {\(\i\)}
    {\(-\i\)}
    {1}
    {\(-1\)}
\end{problem}

\begin{answer}
A
\end{answer}

\begin{solution}
因为 \(607=4\times151+3\)，所以 \(\i^{607}=\i^3=-\i\)，其共轭复数为 \(\i\). 故选 A.
\end{solution}

\begin{problem}
我国古代数学名著《数书九章》有“米谷粒分”题：粮仓开仓收粮，有人送来米 \(1534\) 石，验得米内夹谷，抽样取米一把，数得 \(254\) 粒内夹谷 \(28\) 粒，则这批米内夹谷约为（\quad）

\choices
    {\(134\) 石}
    {\(169\) 石}
    {\(338\) 石}
    {\(1365\) 石}
\end{problem}

\begin{answer}
B
\end{answer}

\begin{solution}
用样本中夹谷粒数所占比例估计整批米中夹谷所占比例，夹谷约有
\(1534\times\dfrac{28}{254}=169.102\cdots\) 石，故约为 \(169\) 石. 故选 B.
\end{solution}

\begin{problem}
已知 \((1+x)^n\) 的展开式中第 \(4\) 项与第 \(8\) 项的二项式系数相等，则奇数项的二项式系数和为（\quad）

\choices
    {\(2^{12}\)}
    {\(2^{11}\)}
    {\(2^{10}\)}
    {\(2^9\)}
\end{problem}

\begin{answer}
D
\end{answer}

\begin{solution}
第 \(4\) 项和第 \(8\) 项的二项式系数分别为 \(\mathrm C_n^3\) 和 \(\mathrm C_n^7\)，故
\(\mathrm C_n^3=\mathrm C_n^7\). 二项式系数关于中间项对称，所以 \(3+7=n\)，从而 \(n=10\).

令 \(x=1\)，有
\(\sum_{k=0}^{10}\mathrm C_{10}^k=2^{10}\). 由于 \(n=10\) 为偶数，奇数项与偶数项的二项式系数和相等，故奇数项的和为 \(2^9\). 故选 D.
\end{solution}

\begin{problem}
设 \(X\sim\mathcal{N}(\mu_1,\sigma_1^2)\)，\(Y\sim\mathcal{N}(\mu_2,\sigma_2^2)\)，这两个正态分布密度曲线如图所示，下列结论中正确的是（\quad）

\bitmapfigure[max width=0.38\linewidth]{img/2015/hubei_science/q4_fig1.png}

\choices
    {\(P(Y\geqslant\mu_2)\geqslant P(Y\geqslant\mu_1)\)}
    {\(P(X\leqslant\sigma_2)\leqslant P(X\leqslant\sigma_1)\)}
    {对任意正数 \(t\)，\(P(X\leqslant t)\geqslant P(Y\leqslant t)\)}
    {对任意正数 \(t\)，\(P(X\geqslant t)\geqslant P(Y\geqslant t)\)}
\end{problem}

\begin{answer}
C
\end{answer}

\begin{solution}
由图可知 \(\mu_1<0<\mu_2\)，且 \(0<\sigma_1<\sigma_2\). 设标准正态分布函数为 \(\Phi\)，则对任意 \(t>0\)，有
\[
\frac{t-\mu_1}{\sigma_1}-\frac{t-\mu_2}{\sigma_2}
=t\left(\frac1{\sigma_1}-\frac1{\sigma_2}\right)+\frac{\mu_2}{\sigma_2}-\frac{\mu_1}{\sigma_1}>0
\]
因此 \(P(X\leqslant t)=\Phi\left(\dfrac{t-\mu_1}{\sigma_1}\right)\geqslant\Phi\left(\dfrac{t-\mu_2}{\sigma_2}\right)=P(Y\leqslant t)\)，选项 C 正确.

又 \(P(Y\geqslant\mu_2)=\dfrac12\)，而 \(\mu_1<\mu_2\) 时 \(P(Y\geqslant\mu_1)>\dfrac12\)，故 A 错. 对任意 \(t>0\)，由标准正态分布函数的单调性，
\(P(X\geqslant t)=1-\Phi\left(\dfrac{t-\mu_1}{\sigma_1}\right)\leqslant1-\Phi\left(\dfrac{t-\mu_2}{\sigma_2}\right)=P(Y\geqslant t)\)，故 D 错. 由于 \(\sigma_2>\sigma_1\)，且 \(X\) 的分布函数严格递增，有 \(P(X\leqslant\sigma_2)>P(X\leqslant\sigma_1)\)，故 B 错，选 C.
\end{solution}

\begin{problem}
设 \(a_1\)，\(a_2\)，\(\cdots\)，\(a_n\in\R\)，\(n\geqslant3\). 若 \(p\)：\(a_1\)，\(a_2\)，\(\cdots\)，\(a_n\) 成等比数列；\(q\)：\((a_1^2+a_2^2+\cdots+a_{n-1}^2)(a_2^2+a_3^2+\cdots+a_n^2)=(a_1a_2+a_2a_3+\cdots+a_{n-1}a_n)^2\)，则（\quad）

\choices
    {\(p\) 是 \(q\) 的充分条件，但不是 \(q\) 的必要条件}
    {\(p\) 是 \(q\) 的必要条件，但不是 \(q\) 的充分条件}
    {\(p\) 是 \(q\) 的充分必要条件}
    {\(p\) 既不是 \(q\) 的充分条件，也不是 \(q\) 的必要条件}
\end{problem}

\begin{answer}
A
\end{answer}

\begin{solution}
若数列成等比数列，设公比为 \(r\)，则 \(a_{k+1}=r a_k\). 于是向量 \((a_2,a_3,\ldots,a_n)=r(a_1,a_2,\ldots,a_{n-1})\)，柯西不等式取等号，故 \(p\) 是 \(q\) 的充分条件.

反之，取 \(a_1=a_2=\cdots=a_{n-1}=0\)，\(a_n=1\)，则 \(q\) 两边均为 \(0\)，但该数列不是等比数列. 因而 \(q\) 不是 \(p\) 的充分条件，即 \(p\) 不是 \(q\) 的必要条件. 故选 A.
\end{solution}

\begin{problem}
已知符号函数 \(\operatorname{sgn}x=\begin{cases}
1, & x>0,\\
0, & x=0,\\
-1, & x<0,
\end{cases}\)，\(f(x)\) 是 \(\R\) 上的增函数，\(g(x)=f(x)-f(ax)\)（\(a>1\)），则（\quad）

\choices
    {\(\operatorname{sgn}[g(x)]=\operatorname{sgn}x\)}
    {\(\operatorname{sgn}[g(x)]=-\operatorname{sgn}x\)}
    {\(\operatorname{sgn}[g(x)]=\operatorname{sgn}[f(x)]\)}
    {\(\operatorname{sgn}[g(x)]=-\operatorname{sgn}[f(x)]\)}
\end{problem}

\begin{answer}
B
\end{answer}

\begin{solution}
当 \(x>0\) 时，\(ax>x\)，所以 \(f(ax)>f(x)\)，从而 \(g(x)<0\). 当 \(x<0\) 时，\(ax<x\)，所以 \(f(ax)<f(x)\)，从而 \(g(x)>0\). 当 \(x=0\) 时，\(g(0)=0\). 因此 \(\operatorname{sgn}[g(x)]=-\operatorname{sgn}x\)，故选 B.
\end{solution}

\begin{problem}
在区间 \([0,1]\) 上随机取两个数 \(x\)，\(y\)，记 \(p_1\) 为事件“\(x+y\geqslant\frac12\)”的概率，\(p_2\) 为事件“\(|x-y|\leqslant\frac12\)”的概率，\(p_3\) 为事件“\(xy\leqslant\frac12\)”的概率，则（\quad）

\choices
    {\(p_1<p_2<p_3\)}
    {\(p_2<p_3<p_1\)}
    {\(p_3<p_1<p_2\)}
    {\(p_3<p_2<p_1\)}
\end{problem}

\begin{answer}
B
\end{answer}

\begin{solution}
在单位正方形 \([0,1]^2\) 中，事件 \(x+y<\frac12\) 是直角等腰三角形，面积为 \(\frac18\)，故 \(p_1=1-\frac18=\frac78\).

事件 \(|x-y|>\frac12\) 是两个边长为 \(\frac12\) 的直角等腰三角形，面积为 \(\frac14\)，故 \(p_2=\frac34\).

事件 \(xy>\frac12\) 的面积为
\(\displaystyle\int_{1/2}^{1}\left(1-\frac{1}{2x}\right)\,\mathrm dx=\frac12-\frac12\ln2\)，故 \(p_3=\frac{1+\ln2}{2}\). 由于 \(\frac12<\ln2<\frac34\)，有
\(p_2=\frac34<p_3<\frac78=p_1\). 其中
\(\ln2=\displaystyle\int_0^1\frac{1}{1+x}\,\mathrm dx\)，且
\(\frac12<\frac{1}{1+x}<1-\frac{x}{2}\)（\(0<x<1\)），所以 \(\frac12<\ln2<\frac34\). 故选 B.
\end{solution}

\begin{problem}
将离心率为 \(e_1\) 的双曲线 \(C_1\) 的实半轴长 \(a\) 和虚半轴长 \(b\)（\(a\neq b\)）同时增加 \(m\)（\(m>0\)）个单位长度，得到离心率为 \(e_2\) 的双曲线 \(C_2\)，则（\quad）

\choices
    {对任意的 \(a,b\)，\(e_1>e_2\)}
    {当 \(a>b\) 时，\(e_1>e_2\)；当 \(a<b\) 时，\(e_1<e_2\)}
    {对任意的 \(a\)，\(b\)，\(e_1<e_2\)}
    {当 \(a>b\) 时，\(e_1<e_2\)；当 \(a<b\) 时，\(e_1>e_2\)}
\end{problem}

\begin{answer}
D
\end{answer}

\begin{solution}
双曲线的离心率为 \(e=\sqrt{1+\dfrac{b^2}{a^2}}\). 比较增加前后的比值，有
\[
\frac{b+m}{a+m}-\frac ba=\frac{m(a-b)}{a(a+m)}
\]
当 \(a>b\) 时，该式为正，所以 \(e_2>e_1\)；当 \(a<b\) 时，该式为负，所以 \(e_2<e_1\). 故选 D.
\end{solution}

\begin{problem}
已知集合 \(A=\{(x,y)\mid x^2+y^2\leqslant1,\ x,y\in\Z\}\)，\(B=\{(x,y)\mid |x|\leqslant2,\ |y|\leqslant2,\ x,y\in\Z\}\)，定义集合 \(A\oplus B=\{(x_1+x_2,y_1+y_2)\mid (x_1,y_1)\in A,\ (x_2,y_2)\in B\}\)，则 \(A\oplus B\) 中元素的个数为（\quad）

\choices
    {77}
    {49}
    {45}
    {30}
\end{problem}

\begin{answer}
C
\end{answer}

\begin{solution}
集合 \(A\) 中的整数点为 \((0,0)\)，\((1,0)\)，\((-1,0)\)，\((0,1)\)，\((0,-1)\). 因而 \(A\oplus B\) 是正方形整数点集 \([-2,2]^2\) 及其向四个坐标方向平移一个单位所得集合的并集.

这些点恰好是 \([-3,3]^2\) 中除四个角点 \((\pm3,\pm3)\) 外的全部整数点：边界上非角点可由相应方向的 \(A\) 中点平移得到，内部点由 \((0,0)\) 得到，而四个角点无法同时在两个坐标方向各越过两单位. 所以元素个数为 \(7\times7-4=45\). 故选 C.
\end{solution}

\begin{problem}
设 \(x\in\R\)，\([x]\) 表示不超过 \(x\) 的最大整数. 若存在实数 \(t\)，使得 \([t]=1\)，\([t^2]=2\)，\(\cdots\)，\([t^n]=n\) 同时成立，则正整数 \(n\) 的最大值是（\quad）

\choices
    {3}
    {4}
    {5}
    {6}
\end{problem}

\begin{answer}
B
\end{answer}

\begin{solution}
条件 \([t^3]=3\) 要求 \(t\geqslant\sqrt[3]{3}\)，而若还满足 \([t^5]=5\)，则 \(t<\sqrt[5]{6}\). 但
\(\sqrt[3]{3}>\sqrt[5]{6}\)，因为两边取十五次方后分别为 \(3^5=243\) 和 \(6^3=216\)，故不可能有 \(n\geqslant5\).

取 \(t=\sqrt[3]{3}\). 因为 \(2^3<3^2<3^3\)，所以 \(1<t<2\) 且 \(2<t^2<3\)；又 \(t^3=3\)，并且
\(4<t^4<5\)，这是因为 \(4^3<3^4<5^3\). 所以 \([t]=1\)，\([t^2]=2\)，\([t^3]=3\)，\([t^4]=4\). 因而最大值为 \(4\). 故选 B.
\end{solution}

\section{填空题}

\begin{problem}
已知向量 \(\overrightarrow{OA}\perp\overrightarrow{AB}\)，\(|\overrightarrow{OA}|=3\)，则 \(\overrightarrow{OA}\cdot\overrightarrow{OB}=\)\fillinblank{}.
\end{problem}

\begin{answer}
9
\end{answer}

\begin{solution}
由 \(\overrightarrow{OB}=\overrightarrow{OA}+\overrightarrow{AB}\)，得
\(\overrightarrow{OA}\cdot\overrightarrow{OB}=\overrightarrow{OA}\cdot\overrightarrow{OA}+\overrightarrow{OA}\cdot\overrightarrow{AB}=|\overrightarrow{OA}|^2=9\).
\end{solution}

\begin{problem}
函数 \(f(x)=4\cos^2\dfrac{x}{2}\cos\left(\dfrac{\pi}{2}-x\right)-2\sin x-|\ln(x+1)|\) 的零点个数为\fillinblank{}.
\end{problem}

\begin{answer}
2
\end{answer}

\begin{solution}
定义域为 \(x>-1\)，且
\[
4\cos^2\frac{x}{2}\cos\left(\frac\pi2-x\right)-2\sin x
=2(1+\cos x)\sin x-2\sin x=\sin2x
\]
所以 \(f(x)=\sin2x-|\ln(x+1)|\).

当 \(-1<x<0\) 时，\(\sin2x<0\)，\(\ln(x+1)<0\)，故 \(f(x)=\sin2x+\ln(x+1)<0\)；而 \(x=0\) 是一个零点.

当 \(0<x\leqslant\frac\pi4\) 时，利用 \(\sin2x\geqslant\frac{4x}{\pi}>x\geqslant\ln(1+x)\)，得 \(f(x)>0\). 在 \([\frac\pi4,\frac\pi2]\) 上，\(\sin2x\) 严格递减而 \(\ln(1+x)\) 严格递增，且两端的函数值一正一负，所以恰有一个零点.

在 \([\frac\pi2,\pi]\) 上，\(\sin2x\leqslant0<\ln(1+x)\)，故无零点. 当 \(x\geqslant\pi\) 且 \(\sin2x\leqslant0\) 时，\(f(x)<0\)；当 \(\sin2x>0\) 时，\(\sin2x\leqslant1<\ln(1+x)\)，其中 \(1+\pi>\e\)，故仍有 \(f(x)<0\). 因此零点共有 \(2\) 个.
\end{solution}

\begin{problem}
如图，一辆汽车在一条水平的公路上向正西行驶，到 \(A\) 处时测得公路北侧一山顶 \(D\) 在西偏北 \(30^\circ\) 的方向上，行驶 \(600\,\mathrm{m}\) 后到达 \(B\) 处，测得此山顶在西偏北 \(75^\circ\) 的方向上，仰角为 \(30^\circ\)，则此山的高度 \(CD=\)\fillinblank{}\(\mathrm{m}\).

\bitmapfigure[max width=0.46\linewidth]{img/2015/hubei_science/q13_fig1.png}
\end{problem}

\begin{answer}
\(100\sqrt{6}\)
\end{answer}

\begin{solution}
设 \(C\) 为山顶 \(D\) 在水平面上的射影，则在水平面三角形 \(ABC\) 中，\(AB=600\,\mathrm m\). 由于 \(AC\) 为西偏北 \(30^\circ\)，而 \(BC\) 为西偏北 \(75^\circ\)，所以 \(\angle A=30^\circ\)，\(\angle B=180^\circ-75^\circ=105^\circ\)，从而 \(\angle C=45^\circ\).

由正弦定理，
\(\displaystyle\frac{BC}{\sin30^\circ}=\frac{AB}{\sin45^\circ}\)，故 \(BC=300\sqrt2\,\mathrm m\). 在直角三角形 \(BCD\) 中，仰角为 \(30^\circ\)，所以
\(CD=BC\tan30^\circ=300\sqrt2\times\dfrac1{\sqrt3}=100\sqrt6\,\mathrm m\).
\end{solution}

\begin{problem}
如图，圆 \(C\) 与 \(x\) 轴相切于点 \(T(1,0)\)，与 \(y\) 轴正半轴交于两点 \(A\)，\(B\)（\(B\) 在 \(A\) 的上方），且 \(|AB|=2\).

\begin{enumerate}
\item 圆 \(C\) 的标准方程为\fillinblank{}；
\item 过点 \(A\) 任作一条直线与圆 \(O\colon x^2+y^2=1\) 相交于 \(M\)，\(N\) 两点，下列三个结论：
\begin{circlelist}
\item \(\dfrac{|NA|}{|NB|}=\dfrac{|MA|}{|MB|}\)；
\item \(\dfrac{|NB|}{|NA|}-\dfrac{|MA|}{|MB|}=2\)；
\item \(\dfrac{|NB|}{|NA|}+\dfrac{|MA|}{|MB|}=2\sqrt2\).
\end{circlelist}
其中正确结论的序号是\fillinblank{}.（写出所有正确结论的序号）
\end{enumerate}

\bitmapfigure[max width=0.30\linewidth]{img/2015/hubei_science/q14_fig1.png}
\end{problem}

\begin{answer}
\((x-1)^2+(y-\sqrt{2})^2=2\)；\(\circled{1}\)，\(\circled{2}\)，\(\circled{3}\)
\end{answer}

\begin{solution}
设圆心为 \(C=(1,k)\). 因为圆与 \(x\) 轴相切于 \(T(1,0)\)，半径为 \(|k|\)，故圆方程为 \((x-1)^2+(y-k)^2=k^2\). 令 \(x=0\)，交点纵坐标满足
\(1+y^2-2ky=0\). 两根之差为 \(2\sqrt{k^2-1}=2\)，且两交点在 \(y\) 轴正半轴上，所以 \(k=\sqrt2\). 因而圆的标准方程为 \((x-1)^2+(y-\sqrt2)^2=2\)，并且
\(A=(0,\sqrt2-1)\)，\(B=(0,\sqrt2+1)\).

记 \(a=\sqrt2-1\)，\(b=\sqrt2+1\)，则 \(ab=1\). 对单位圆上任意一点 \(U=(u,v)\)，有

\(UA^2=u^2+(v-a)^2=1+a^2-2av\)，而
\(a^2UB^2=a^2\bigl(u^2+(v-b)^2\bigr)=a^2+1-2av=UA^2\).

由于距离均为正数，\(\dfrac{UA}{UB}=a=\sqrt2-1\). 分别取 \(U=M\)，\(N\)，得到
\(\dfrac{|MA|}{|MB|}=\dfrac{|NA|}{|NB|}=\sqrt2-1\).

于是 \(\dfrac{|NB|}{|NA|}=\sqrt2+1\)，从而
\(\dfrac{|NA|}{|NB|}=\dfrac{|MA|}{|MB|}\)，
\(\dfrac{|NB|}{|NA|}-\dfrac{|MA|}{|MB|}=2\)，以及
\(\dfrac{|NB|}{|NA|}+\dfrac{|MA|}{|MB|}=2\sqrt2\). 三个结论均正确.
\end{solution}

\section{选考题：填空题}

\begin{problem}
如图，\(PA\) 是圆的切线，\(A\) 为切点，\(PBC\) 是圆的割线，且 \(BC=3PB\)，则 \(\dfrac{AB}{AC}=\)\fillinblank{}.
\end{problem}

\begin{answer}
\(\dfrac{1}{2}\)
\end{answer}

\begin{solution}
由切割线定理，\(PA^2=PB\cdot PC=PB(PB+BC)=4PB^2\)，故 \(PA=2PB\). 由切线与弦所成的角等于弦所对的圆周角，且 \(PB\) 与 \(PC\) 是同一条射线，三角形 \(PAB\) 与三角形 \(PCA\) 相似，其中 \(PB\leftrightarrow PA\)，\(AB\leftrightarrow AC\)，因此
\(\dfrac{AB}{AC}=\dfrac{PB}{PA}=\dfrac12\).
\end{solution}

\begin{problem}
在直角坐标系 \(xOy\) 中，以 \(O\) 为极点，\(x\) 轴的正半轴为极轴建立极坐标系. 已知直线 \(l\) 的极坐标方程为 \(\rho(\sin\theta-3\cos\theta)=0\)，曲线 \(C\) 的参数方程为 \(\begin{cases}
x=t-\dfrac1t,\\
y=t+\dfrac1t
\end{cases}\)（\(t\) 为参数），\(l\) 与 \(C\) 相交于 \(A\)，\(B\) 两点，则 \(|AB|=\)\fillinblank{}.
\end{problem}

\begin{answer}
\(2\sqrt{5}\)
\end{answer}

\begin{solution}
由 \(x=\rho\cos\theta\)，\(y=\rho\sin\theta\)，直线方程化为 \(y-3x=0\). 代入参数方程，得
\(t+\dfrac1t=3\left(t-\dfrac1t\right)\)，即 \(t^2=2\)，所以 \(t=\pm\sqrt2\).

对应交点为 \(\left(\dfrac1{\sqrt2},\dfrac3{\sqrt2}\right)\) 与 \(\left(-\dfrac1{\sqrt2},-\dfrac3{\sqrt2}\right)\)，故
\[
|AB|=2\sqrt{\left(\frac1{\sqrt2}\right)^2+\left(\frac3{\sqrt2}\right)^2}=2\sqrt5
\]
\end{solution}

\section{解答题}

\begin{problem}
某同学用“五点法”画函数 \(f(x)=A\sin(\omega x+\varphi)\)（\(\omega>0\)，\(|\varphi|<\dfrac{\pi}{2}\)）在某一个周期内的图象时，列表并填入了部分数据，如下表：

\begin{center}
\renewcommand{\arraystretch}{1.25}
\begin{tabular}{|c|c|c|c|c|c|}
\hline
\(\omega x+\varphi\) & 0 & \(\dfrac{\pi}{2}\) & \(\pi\) & \(\dfrac{3\pi}{2}\) & \(2\pi\) \\
\hline
\(x\) &  & \(\dfrac{\pi}{3}\) &  & \(\dfrac{5\pi}{6}\) &  \\
\hline
\(A\sin(\omega x+\varphi)\) & 0 & 5 &  & \(-5\) & 0 \\
\hline
\end{tabular}
\end{center}

\begin{enumerate}
\item 请将上表数据补充完整，并直接写出函数 \(f(x)\) 的解析式；
\item 将 \(y=f(x)\) 图象上所有点向左平行移动 \(\theta\)（\(\theta>0\)）个单位长度，得到 \(y=g(x)\) 的图象. 若 \(y=g(x)\) 图象的一个对称中心为 \(\left(\dfrac{5\pi}{12},0\right)\)，求 \(\theta\) 的最小值.
\end{enumerate}
\end{problem}

\begin{solution}
\begin{enumerate}
\item 由表中两列数据，\(\omega\left(\dfrac{5\pi}{6}-\dfrac{\pi}{3}\right)=\pi\)，所以 \(\omega=2\). 再由 \(2\cdot\dfrac{\pi}{3}+\varphi=\dfrac\pi2\)，得 \(\varphi=-\dfrac\pi6\). 由 \(A\sin\dfrac\pi2=5\) 得 \(A=5\).

当相位分别为 \(0\)，\(\dfrac\pi2\)，\(\pi\)，\(\dfrac{3\pi}{2}\)，\(2\pi\) 时，\(x\) 分别为
\(\dfrac\pi{12}\)，\(\dfrac\pi3\)，\(\dfrac{7\pi}{12}\)，\(\dfrac{5\pi}{6}\)，\(\dfrac{13\pi}{12}\)，函数值分别为 \(0\)，\(5\)，\(0\)，\(-5\)，\(0\). 因而
\(f(x)=5\sin\left(2x-\dfrac\pi6\right)\).

\item 向左平移 \(\theta\) 后，\(g(x)=f(x+\theta)=5\sin\left(2x+2\theta-\dfrac\pi6\right)\). 正弦函数图象的对称中心满足相位为 \(k\pi\)，故
\(\frac{5\pi}{12}=\frac\pi{12}-\theta+\frac{k\pi}{2}\)，\(k\in\Z\).
于是 \(\theta=-\dfrac\pi3+\dfrac{k\pi}{2}\). 取使 \(\theta>0\) 的最小整数 \(k=1\)，得 \(\theta_{\min}=\dfrac\pi6\).
\end{enumerate}
\end{solution}

\begin{problem}
设等差数列 \(\{a_n\}\) 的公差为 \(d\)，前 \(n\) 项和为 \(S_n\)，等比数列 \(\{b_n\}\) 的公比为 \(q\). 已知 \(b_1=a_1\)，\(b_2=2\)，\(q=d\)，\(S_{10}=100\).

\begin{enumerate}
\item 求数列 \(\{a_n\}\)，\(\{b_n\}\) 的通项公式；
\item 当 \(d>1\) 时，记 \(c_n=\dfrac{a_n}{b_n}\)，求数列 \(\{c_n\}\) 的前 \(n\) 项和 \(T_n\).
\end{enumerate}
\end{problem}

\begin{solution}
\begin{enumerate}
\item 因为 \(b_2=b_1q=2\)，且 \(b_1=a_1\)，所以 \(a_1=\dfrac2d\). 由 \(S_{10}=5(2a_1+9d)=100\)，得
\(2a_1+9d=20\)，\(\frac4d+9d=20\)
即 \(9d^2-20d+4=0\). 解得 \(d=2\) 或 \(d=\dfrac29\).

当 \(d=2\) 时，\(a_1=1\)，\(b_1=1\)，故 \(a_n=2n-1\)，\(b_n=2^{n-1}\). 当 \(d=\dfrac29\) 时，\(a_1=9\)，\(b_1=9\)，故 \(a_n=\dfrac{2n+79}{9}\)，\(b_n=9\left(\dfrac29\right)^{n-1}\).

\item 由 \(d>1\)，取 \(d=2\)，所以 \(c_n=\dfrac{2n-1}{2^{n-1}}\). 设
\(T_n=6-\dfrac{2n+3}{2^{n-1}}\). 当 \(n=1\) 时右边为 \(1=c_1\)；若对 \(n-1\) 成立，则
\[
T_{n-1}+c_n=6-\frac{2n+1}{2^{n-2}}+\frac{2n-1}{2^{n-1}}
=6-\frac{2n+3}{2^{n-1}}
\]
所以 \(\displaystyle T_n=6-\frac{2n+3}{2^{n-1}}\).
\end{enumerate}
\end{solution}

\begin{problem}
《九章算术》中，将底面为长方形且有一条侧棱与底面垂直的四棱锥称之为阳马，将四个面都为直角三角形的四面体称之为鳖臑. 如图，在阳马 \(P-ABCD\) 中，侧棱 \(PD\perp\) 底面 \(ABCD\)，且 \(PD=CD\)，过棱 \(PC\) 的中点 \(E\)，作 \(EF\perp PB\) 交 \(PB\) 于点 \(F\)，连接 \(DE\)，\(DF\)，\(BD\)，\(BE\).

\begin{enumerate}
\item 证明：\(PB\perp\) 平面 \(DEF\). 试判断四面体 \(DBEF\) 是否为鳖臑，若是，写出其每个面的直角（只需写出结论）；若不是，说明理由；
\item 若面 \(DEF\) 与面 \(ABCD\) 所成二面角的大小为 \(\dfrac\pi3\)，求 \(\dfrac{DC}{BC}\) 的值.
\end{enumerate}

\bitmapfigure[max width=0.44\linewidth]{img/2015/hubei_science/q19_fig1.png}
\end{problem}

\begin{solution}
\begin{enumerate}
\item
设 \(DC=a\)，\(BC=b\)，建立空间直角坐标系，使
\(D=(0,0,0)\)，\(C=(a,0,0)\)，\(B=(a,b,0)\)，\(A=(0,b,0)\)，\(P=(0,0,a)\). 则
\(E=\left(\dfrac a2,0,\dfrac a2\right)\)，且 \(\overrightarrow{PB}=(a,b,-a)\).

因为 \(\overrightarrow{DE}=\left(\dfrac a2,0,\dfrac a2\right)\)，所以 \(\overrightarrow{DE}\cdot\overrightarrow{PB}=0\). 又 \(EF\perp PB\)，且 \(DE\) 与 \(EF\) 是平面 \(DEF\) 内相交的两条直线，故 \(PB\perp\) 平面 \(DEF\).

设 \(F=P+\lambda\overrightarrow{PB}\). 由 \(EF\perp PB\)，得 \(\lambda=\dfrac{a^2}{2a^2+b^2}\)，从而
\(F=(\lambda a,\lambda b,a(1-\lambda))\)，\(\overrightarrow{EF}=\left(a\left(\frac12-\lambda\right),-\lambda b,a\left(\lambda-\frac12\right)\right)\).
\(\overrightarrow{DE}\cdot\overrightarrow{EB} =\frac{a^2}{4}-\frac{a^2}{4}=0\)，\(\overrightarrow{DE}\cdot\overrightarrow{EF} =\frac{a^2}{2}\left(\frac12-\lambda\right)+\frac{a^2}{2}\left(\lambda-\frac12\right)=0\).
又因 \(EF\perp PB\) 且 \(FB\parallel PB\)，有 \(\overrightarrow{EF}\cdot\overrightarrow{FB}=0\). 最后
\[
\overrightarrow{DF}\cdot\overrightarrow{FB}
=(1-\lambda)\left(\lambda(2a^2+b^2)-a^2\right)=0
\]
因此四面体 \(DBEF\) 的四个面均为直角三角形，是鳖臑，四个面的直角分别为 \(\angle DEB\)，\(\angle DEF\)，\(\angle EFB\)，\(\angle DFB\).

\item 平面 \(DEF\) 的一个法向量可取 \(\overrightarrow{PB}=(a,b,-a)\)，底面 \(ABCD\) 的法向量可取 \((0,0,1)\). 因此两平面所成的锐二面角 \(\alpha\) 满足
\[
\cos\alpha=\frac{a}{\sqrt{2a^2+b^2}}
\]
由 \(\alpha=\dfrac\pi3\)，得 \(\dfrac a{\sqrt{2a^2+b^2}}=\dfrac12\)，即 \(b^2=2a^2\). 所以
\(\displaystyle\frac{DC}{BC}=\frac ab=\frac{\sqrt2}{2}\).
\end{enumerate}
\end{solution}

\begin{problem}
某厂用鲜牛奶在某台设备上生产 A，B 两种奶制品. 生产 \(1\) 吨 A 产品需鲜牛奶 \(2\) 吨，使用设备 \(1\) 小时，获利 \(1000\) 元；生产 \(1\) 吨 B 产品需鲜牛奶 \(1.5\) 吨，使用设备 \(1.5\) 小时，获利 \(1200\) 元. 要求每天 B 产品的产量不超过 A 产品产量的 \(2\) 倍，设备每天生产 A，B 两种产品时间之和不超过 \(12\) 小时. 假定每天可获取的鲜牛奶数量 \(W\)（单位：吨）是一个随机变量，其分布列为

\begin{center}
\renewcommand{\arraystretch}{1.2}
\begin{tabular}{|c|c|c|c|}
\hline
\(W\) & 12 & 15 & 18 \\
\hline
\(P\) & 0.3 & 0.5 & 0.2 \\
\hline
\end{tabular}
\end{center}

该厂每天根据获取的鲜牛奶数量安排生产，使其获利最大，因此每天的最大获利 \(Z\)（单位：元）是一个随机变量.

\begin{enumerate}
\item 求 \(Z\) 的分布列和均值；
\item 若每天可获取的鲜牛奶数量相互独立，求 \(3\) 天中至少有 \(1\) 天的最大获利超过 \(10000\) 元的概率.
\end{enumerate}
\end{problem}

\begin{solution}
\begin{enumerate}
\item
设每天生产 A，B 两种产品的产量分别为 \(x\)，\(y\)，则
\(x\geqslant0\)，\(y\geqslant0\)，\(y\leqslant2x\)，\(2x+1.5y\leqslant W\)，\(x+1.5y\leqslant12\).
利润为 \(1000x+1200y\).

当 \(W=12\) 时，可行域顶点为 \((0,0)\)，\((6,0)\)，\(\left(\dfrac{12}{5},\dfrac{24}{5}\right)\)，对应利润分别为 \(0\)，\(6000\)，\(1000\times\dfrac{12}{5}+1200\times\dfrac{24}{5}=8160\)，故最优点为 \(\left(\dfrac{12}{5},\dfrac{24}{5}\right)\)；当 \(W=15\) 时，可行域顶点为 \((0,0)\)，\(\left(\dfrac{15}{2},0\right)\)，\((3,6)\)，对应利润分别为 \(0\)，\(7500\)，\(10200\)，故最优点为 \((3,6)\)；当 \(W=18\) 时，可行域顶点为 \((0,0)\)，\((9,0)\)，\((6,4)\)，\((3,6)\)，对应利润分别为 \(0\)，\(9000\)，\(10800\)，\(10200\) 元，故最优点为 \((6,4)\).

因此 \(Z\) 的分布列为
\begin{center}
\begin{tabular}{c|ccc}
\(Z\)&\(8160\)&\(10200\)&\(10800\)\\\hline
\(P\)&\(0.3\)&\(0.5\)&\(0.2\)
\end{tabular}
\end{center}
且
\(E(Z)=8160\times0.3+10200\times0.5+10800\times0.2=9708\).

\item \(Z>10000\) 对应 \(W=15\) 或 \(W=18\)，其概率为 \(0.5+0.2=0.7\). 由三天相互独立，所求概率为
\(1-(1-0.7)^3=1-0.3^3=0.973\).
\end{enumerate}
\end{solution}

\begin{problem}
一种作图工具如图 1 所示. \(O\) 是滑槽 \(AB\) 的中点，短杆 \(ON\) 可绕 \(O\) 转动，长杆 \(MN\) 通过 \(N\) 处铰链与 \(ON\) 连接，\(MN\) 上的栓子 \(D\) 可沿滑槽 \(AB\) 滑动，且 \(DN=ON=1\)，\(MN=3\). 当栓子 \(D\) 在滑槽 \(AB\) 内作往复运动时，带动 \(N\) 绕 \(O\) 转动一周（\(D\) 不动时，\(N\) 也不动），\(M\) 处的笔尖画出的曲线记为 \(C\). 以 \(O\) 为原点，\(AB\) 所在的直线为 \(x\) 轴建立如图 2 所示的平面直角坐标系.

\begin{enumerate}
\item 求曲线 \(C\) 的方程；
\item 设动直线 \(l\) 与两定直线 \(l_1\colon x-2y=0\) 和 \(l_2\colon x+2y=0\) 分别交于 \(P\)，\(Q\) 两点. 若直线 \(l\) 总与曲线 \(C\) 有且只有一个公共点，试探究：\(\triangle OPQ\) 的面积是否存在最小值？若存在，求出该最小值；若不存在，说明理由.
\end{enumerate}

\FigureLayoutDeclare{img/2015/hubei_science/q21_fig1.png}{17.0cm}{40bp 10bp 40bp 14bp}
\bitmapfigure[max width=0.68\linewidth]{img/2015/hubei_science/q21_fig1.png}
\end{problem}

\begin{solution}
\begin{enumerate}
\item 设 \(N=(\cos t,\sin t)\). 由 \(DN=ON=1\) 且 \(D\) 在 \(x\) 轴上，除去静止的分支 \(D=O\) 外，运动分支满足 \(D=(2\cos t,0)\). 因为 \(D\) 在 \(MN\) 上且 \(DN=1\)，\(M\) 是从 \(N\) 经过 \(D\) 延长三倍距离的点，所以
\[
M= N+3(D-N)=(4\cos t,-2\sin t)
\]
故曲线 \(C\) 的方程为 \(\displaystyle\frac{x^2}{16}+\frac{y^2}{4}=1\).

\item 直线 \(l\) 与椭圆只有一个公共点，所以它是椭圆的切线. 设切点为 \((4\cos t,2\sin t)\)，切线方程为
\[
x\cos t+2y\sin t=4
\]
令 \(P=(p,\frac p2)\)，\(Q=(q,-\frac q2)\). 在题设两交点存在时，\(\cos t+\sin t\) 与 \(\cos t-\sin t\) 均不为零. 将切线分别与 \(l_1\)，\(l_2\) 联立，得
\(p=\dfrac4{\cos t+\sin t}\)，\(q=\dfrac4{\cos t-\sin t}\). 因而
\[
S_{\triangle OPQ}=\frac12\left|\det\begin{pmatrix}p&p/2\\q&-q/2\end{pmatrix}\right|
=\frac{|pq|}{2}
=\frac8{|\cos^2t-\sin^2t|}
=\frac8{|\cos2t|}\geqslant8
\]
当 \(t=0\) 或 \(t=\dfrac\pi2\) 时取等号，故面积存在最小值，且最小值为 \(8\).
\end{enumerate}
\end{solution}

\begin{problem}
已知数列 \(\{a_n\}\) 的各项均为正数，\(b_n=n\left(1+\dfrac1n\right)^n a_n\)（\(n\in\mathbf{N}_+\)），\(\e\) 为自然对数的底数.

\begin{enumerate}
\item 求函数 \(f(x)=1+x-\e^x\) 的单调区间，并比较 \(\left(1+\dfrac1n\right)^n\) 与 \(\e\) 的大小；
\item 计算 \(\dfrac{b_1}{a_1}\)，\(\dfrac{b_1b_2}{a_1a_2}\)，\(\dfrac{b_1b_2b_3}{a_1a_2a_3}\)，由此推测计算 \(\dfrac{b_1b_2\cdots b_n}{a_1a_2\cdots a_n}\) 的公式，并给出证明；
\item 令 \(c_n=(a_1a_2\cdots a_n)^{\frac1n}\)，数列 \(\{a_n\}\)，\(\{c_n\}\) 的前 \(n\) 项和分别记为 \(S_n\)，\(T_n\)，证明：\(T_n<\e S_n\).
\end{enumerate}
\end{problem}

\begin{solution}
\begin{enumerate}
\item \(f'(x)=1-\e^x\). 当 \(x<0\) 时 \(f'(x)>0\)，当 \(x>0\) 时 \(f'(x)<0\)，所以 \(f(x)\) 在 \((-\infty,0)\) 上单调递增，在 \((0,+\infty)\) 上单调递减. 又 \(f(0)=0\)，故对 \(n\in\mathbf N_+\)，由 \(f(1/n)<f(0)\) 得 \(1+\dfrac1n<\e^{1/n}\)，从而 \(\left(1+\dfrac1n\right)^n<\e\).

\item 由定义，
\(\frac{b_1}{a_1}=2\)，\(\frac{b_1b_2}{a_1a_2}=2\times2\left(\frac32\right)^2=9\)
\[
\frac{b_1b_2b_3}{a_1a_2a_3}=2\times2\left(\frac32\right)^2\times3\left(\frac43\right)^3=64
\]
一般地，
\[
\frac{b_1b_2\cdots b_n}{a_1a_2\cdots a_n}
=\prod_{k=1}^n k\left(\frac{k+1}{k}\right)^k=(n+1)^n
\]
其中最后一个等式由约去相邻幂次因子得到. 也可写成
\(\prod_{k=1}^n k\cdot\dfrac{\prod_{k=1}^n(k+1)^k}{\prod_{k=1}^n k^k}=(n+1)^n\)，故公式成立.

\item 记 \(\lambda_k=k\left(1+\dfrac1k\right)^k=\dfrac{b_k}{a_k}\). 由第（2）问，\(\prod_{k=1}^j\lambda_k=(j+1)^j\). 因而
\[
c_j=\frac{\left(\prod_{k=1}^j\lambda_k a_k\right)^{1/j}}{j+1}
\leqslant\frac{1}{j(j+1)}\sum_{k=1}^j\lambda_k a_k
\]
其中不等号由算术平均数不小于几何平均数得到. 对 \(j=1\)，\(2\)，\(\ldots\)，\(n\) 求和并交换求和次序，得
\[
T_n\leqslant\sum_{k=1}^n\lambda_k a_k\sum_{j=k}^n\frac1{j(j+1)}
=\sum_{k=1}^n\lambda_k a_k\left(\frac1k-\frac1{n+1}\right)
\]
由于 \(a_k>0\)，且
\[
\lambda_k\left(\frac1k-\frac1{n+1}\right)<\frac{\lambda_k}{k}=\left(1+\frac1k\right)^k<\e
\]
所以
\(T_n<\e\sum_{k=1}^n a_k=\e S_n\).
\end{enumerate}
\end{solution}
